%% kapa: exam=1 round=36 year=2025 period=2 form=A subject=appraisal_law subject_ko=감정평가관계법규 seq=36 number=36
%% source: Q-Net 감정평가사 기출문제 (cst00309) artlSeq=5231525 file=raw/1008/5231525/2025년 제36회 감정평가사 1차시험 2교시 문제지 원본.pdf
%% conversion: kordoc 4.24.51 -> tools/gen1.py (auto); review=proofread-sonnet; flags=none
\begin{question}{36}
동산ㆍ채권 등의 담보에 관한 법령상 동산담보권에 관한 설명으로 옳은 것은?
\begin{choices}
  \item 담보권설정자의 사업자등록이 말소된 경우에는 이미 설정된 동산담보권의 효력에도 영향을 미친다.
  \item 담보권자는 공탁금의 회수를 청구할 수 없다.
  \item 동산담보권의 효력은 담보목적물에 대한 압류 후에 담보권설정자가 그 담보목적물로부터 수취한 과실(果實)에 미치지 않는다.
  \item 동산담보권은 피담보채권과 분리하여 타인에게 양도할 수 있다.
  \item 담보권자는 동산담보권을 방해하는 자에게 방해의 제거를 청구할 수 있지만, 동산담보권을 방해할 우려가 있다는 이유로 방해의 예방을 청구할 수는 없다.
\end{choices}
\end{question}
\answer{2}
